\documentclass[11pt,a4paper]{article}

% =========================================================================
% PAKET-PAKET UTAMA & PENGATURAN TATA LETAK
% =========================================================================
\usepackage[utf8]{inputenc}
\usepackage[T1]{fontenc}
\usepackage{lmodern}
\usepackage[indonesian]{babel}
\usepackage{amsmath,amssymb,amsfonts,amsthm}
\usepackage{bm}
\usepackage{geometry}
\geometry{top=1.8cm, bottom=1.8cm, left=1.8cm, right=1.8cm}
\usepackage{graphicx}
\usepackage{booktabs}
\usepackage{tabularx}
\usepackage{float}
\usepackage{microtype}
\usepackage{xcolor}
\usepackage{mdframed}
\usepackage{caption}
\usepackage{subcaption}
\usepackage{hyperref}

\hypersetup{
    colorlinks=true,
    linkcolor=teal!70!black,
    citecolor=teal!70!black,
    urlcolor=teal!70!black,
    pdfencoding=auto,
    psdextra
}

% Reduksi spasi vertikal persamaan
\setlength{\abovedisplayskip}{6pt}
\setlength{\belowdisplayskip}{6pt}
\setlength{\abovedisplayshortskip}{3pt}
\setlength{\belowdisplayshortskip}{3pt}

% Pengaturan warna
\definecolor{dgpblue}{RGB}{26, 82, 118}
\definecolor{dgpgreen}{RGB}{30, 132, 73}
\definecolor{dgppurple}{RGB}{108, 52, 131}
\definecolor{dgpred}{RGB}{148, 49, 38}
\definecolor{boxbg}{RGB}{245, 248, 250}
\definecolor{codebg}{RGB}{240, 242, 245}

% Kotak informasi mdframed
\newmdenv[
    linecolor=dgpblue,
    outerlinewidth=1.2pt,
    roundcorner=4pt,
    innertopmargin=7pt,
    innerbottommargin=7pt,
    innerrightmargin=10pt,
    innerleftmargin=10pt,
    backgroundcolor=boxbg
]{mathbox}

\newmdenv[
    linecolor=dgpred,
    outerlinewidth=1.0pt,
    roundcorner=3pt,
    innertopmargin=6pt,
    innerbottommargin=6pt,
    innerrightmargin=8pt,
    innerleftmargin=8pt,
    backgroundcolor=red!3
]{alertbox}

\newmdenv[
    linecolor=gray!60,
    outerlinewidth=0.8pt,
    roundcorner=3pt,
    innertopmargin=5pt,
    innerbottommargin=5pt,
    innerrightmargin=8pt,
    innerleftmargin=8pt,
    backgroundcolor=codebg
]{codebox}

% =========================================================================
% JUDUL DOKUMEN
% =========================================================================
\title{
    \vspace{-1.2cm}
    \textbf{\Large SUPLEMEN KAJIAN KRITIS: NOTEBOOK \texttt{CO2\_1.ipynb}} \\[0.3em]
    \textbf{\large Evaluasi Algoritma: Status DSVI, ELBO, dan Optimasi Hyperparameter}
}
\author{
    \textbf{Topik Tugas Akhir}: \textit{Image Segmentation with Deep Gaussian Process} \\
    Catatan Pendalaman Teori \& Implementasi Komputasi Bimbingan 06
}
\date{\today}

\begin{document}

\maketitle
\vspace{-0.5cm}

\begin{abstract}
\noindent Dokumen teknis ini disusun untuk menjawab secara analitik, terperinci, dan berbasis bukti kode sumber tiga pertanyaan mendasar terkait notebook eksperimen \texttt{CO2\_1.ipynb}: (1) Apakah DGP pada notebook tersebut menggunakan \textit{Doubly Stochastic Variational Inference} (DSVI)? (2) Apakah terdapat perhitungan dan pemanfaatan \textit{Evidence Lower Bound} (ELBO)? dan (3) Apakah dilakukan prosedur optimasi hyperparameter? Berdasarkan bedah kode dan analisis struktur probabilitas, ditemukan bahwa model yang berjalan adalah \textbf{Kaskade Gaussian Process Eksak Independen (\textit{Cascaded ExactGP})} dengan parameter tetap (\textit{hardcoded}), \textbf{BUKAN} model \textit{Deep Gaussian Process} sejati berbasis DSVI dan ELBO. Dokumen ini juga menyajikan tinjauan matematis formal bagi mahasiswa Departemen Matematika dalam mempertahankan integritas teori pada draf Tugas Akhir.
\end{abstract}

\vspace{0.1cm}
\hrule
\vspace{0.3cm}

% =========================================================================
% RINGKASAN EKSEKUTIF
% =========================================================================
\section{Ringkasan Eksekutif: Jawaban 3 Pertanyaan Kunci}
\label{sec:ringkasan}

\begin{mathbox}
\begin{center}
\textbf{\large TABEL KESIMPULAN HASIL AUDIT KODE \texttt{CO2\_1.ipynb}}
\end{center}
\vspace{0.1cm}
\begin{tabularx}{\textwidth}{p{4.2cm} c X}
\toprule
\textbf{Pertanyaan} & \textbf{Status} & \textbf{Keterangan Matematis \& Implementasi} \\
\midrule
\textbf{1. Menggunakan DSVI?} & \textcolor{dgpred}{\textbf{TIDAK}} & Menggunakan kelas \texttt{gpytorch.models.ExactGP} berdimensi $[N \times N]$ penuh ($N = 1.289$). Tidak ada \textit{inducing points} $\mathbf{Z}$, tidak ada distribusi variasional $q(\mathbf{U})$, dan tidak ada inferensi gabungan antar-lapisan. \\
\addlinespace[4pt]
\textbf{2. Menggunakan ELBO?} & \textcolor{dgpred}{\textbf{TIDAK}} & Tidak terdapat objek loss \texttt{VariationalELBO} maupun \texttt{ExactMarginalLogLikelihood}. Tidak ada \textit{training loop} dan optimasi gradien. Model langsung diposisikan pada mode evaluasi prediktif (\texttt{eval()}). \\
\addlinespace[4pt]
\textbf{3. Melakukan Optimasi Hyperparameter?} & \textcolor{dgpred}{\textbf{TIDAK}} & Seluruh hyperparameter kernel ($\ell$, $\sigma_f^2$, $\sigma_n^2$) dimasukkan secara manual dan bernilai tetap (\textit{fixed/hardcoded values}). Tidak ada algoritma \textit{Gradient Descent}, Adam, maupun L-BFGS. \\
\bottomrule
\end{tabularx}
\end{mathbox}

\vspace{0.2cm}

% =========================================================================
% SEKSI 1: ANALISIS DSVI
% =========================================================================
\section{Kajian Pertanyaan 1: Mengapa Model BUKAN Merupakan DSVI?}
\label{sec:analisis_dsvi}

\noindent Kerangka kerja \textbf{Doubly Stochastic Variational Inference (DSVI)} yang dirumuskan oleh Salimbeni \& Deisenroth (NeurIPS 2017) merupakan paradigma standar bagi DGP modern. Untuk menguji apakah \texttt{CO2\_1.ipynb} menggunakan DSVI, mari kita bandingkan syarat formal DSVI dengan apa yang tertulis pada kode:

\subsection{Karakteristik Mutlak DSVI Sejati (Salimbeni \& Deisenroth, 2017)}
Model DGP $L$-lapisan dengan DSVI menuntut struktur hierarkis probabilitas bersama (\textit{joint distribution}):
\begin{equation}
    p\left(\mathbf{y}, \{\mathbf{f}^{(\ell)}, \mathbf{u}^{(\ell)}\}_{\ell=1}^L\right) = p(\mathbf{y} \mid \mathbf{f}^{(L)}) \prod_{\ell=1}^L p\left(\mathbf{f}^{(\ell)} \mid \mathbf{u}^{(\ell)}, \mathbf{f}^{(\ell-1)}\right) p(\mathbf{u}^{(\ell)})
\end{equation}
dengan aproksimasi distribusi variasional bebas:
\begin{equation}
    q\left(\{\mathbf{f}^{(\ell)}, \mathbf{u}^{(\ell)}\}_{\ell=1}^L\right) = \prod_{\ell=1}^L p\left(\mathbf{f}^{(\ell)} \mid \mathbf{u}^{(\ell)}, \mathbf{f}^{(\ell-1)}\right) q(\mathbf{u}^{(\ell)}), \quad q(\mathbf{u}^{(\ell)}) = \mathcal{N}(\mathbf{m}^{(\ell)}, \mathbf{S}^{(\ell)})
\end{equation}
Inti matematis DSVI mensyaratkan:
\begin{enumerate}\setlength{\itemsep}{1pt}
    \item \textbf{Adanya Titik Induksi (\textit{Inducing Points}) $\mathbf{Z}^{(\ell)} \in \mathbb{R}^{M_\ell \times D_{\ell-1}}$}: Jumlah titik $M \ll N$ digunakan agar komputasi matriks tereduksi menjadi $\mathcal{O}(M^3)$.
    \item \textbf{Trik Reparameterisasi Antar-Lapisan}: Variabel laten pada setiap titik sampel ditransmisikan menggunakan deviasi standar kondisional: $h_{n,s}^{(\ell)} = \mu_{n,s}^{(\ell)} + \epsilon_{n,s} \cdot \sqrt{v_{n,s}^{(\ell)}}$ dengan $\epsilon \sim \mathcal{N}(0, I)$.
    \item \textbf{Inferensi Bersama (\textit{Joint Learning})}: Parameter variasional $\{\mathbf{m}^{(\ell)}, \mathbf{S}^{(\ell)}\}$ pada seluruh lapisan dioptimasi secara serentak dari ujung ke ujung (\textit{end-to-end}).
\end{enumerate}

\subsection{Fakta Implementasi pada Notebook \texttt{CO2\_1.ipynb}}
Ketika menelaah baris kode pada Sel 8, 10, dan 12 pada notebook \texttt{CO2\_1.ipynb}, kita menemukan implementasi berikut:

\begin{codebox}
\small
\textbf{Cuplikan Sel 8: Definisi Kelas Model}
\begin{verbatim}
class GPyTorchExactGP(gpytorch.models.ExactGP):
    def __init__(self, train_x, train_y, likelihood, lengthscale=0.20):
        super().__init__(train_x, train_y, likelihood)
        self.mean_module = gpytorch.means.ConstantMean()
        self.covar_module = gpytorch.kernels.ScaleKernel(gpytorch.kernels.RBFKernel())
        self.covar_module.base_kernel.lengthscale = torch.tensor([[lengthscale]])
\end{verbatim}
\end{codebox}

\begin{codebox}
\small
\textbf{Cuplikan Sel 10: Setup DGP 2-Layer}
\begin{verbatim}
# Membuat target tiruan (dummy) untuk Layer 1:
h1_train = (train_x - train_x.mean()) / train_x.std()

gp1 = GPyTorchExactGP(train_x, h1_train, lh1, lengthscale=0.30)
gp2 = GPyTorchExactGP(h1_train, train_y, lh2, lengthscale=0.30)
\end{verbatim}
\end{codebox}

\noindent Dari konstruksi kode di atas, terlihat kejanggalan teoretis yang sangat mendasar:
\begin{enumerate}
    \item \textbf{Mewarisi \texttt{ExactGP}, Bukan \texttt{ApproximateGP}}: DSVI di dalam pustaka GPyTorch selalu dibangun di atas modul variasional \texttt{gpytorch.models.ApproximateGP} bersama \texttt{VariationalStrategy}. Penggunaan \texttt{ExactGP} membuktikan bahwa model menyelesaikan invers analitik penuh $[K_{NN} + \sigma_n^2 I]^{-1}$ berukuran $1.289 \times 1.289$, bukan aproksimasi variasional berbasis inducing points.
    \item \textbf{Target Laten Lapisan 1 Bersifat Deterministik dan Artifisial}: Variabel $h^{(1)}$ pada DGP sejati adalah variabel acak laten tak teramati yang distribusinya dicari via posterior. Namun pada \texttt{CO2\_1.ipynb}, target lapisan 1 (\texttt{h1\_train}) didefinisikan secara paksa sebagai normalisasi linier dari input waktu $x$:
    \begin{equation}
        h_{\mathrm{train}}^{(1)} = \frac{x_{\mathrm{train}} - \bar{x}}{s_x}
    \end{equation}
    Ini berarti \texttt{gp1} dan \texttt{gp2} adalah \textbf{dua model regresi independen yang terisolasi total}.
    \item \textbf{Status Model Sesungguhnya}: Model ini di dalam literatur tepatnya disebut sebagai \textbf{\textit{Cascaded ExactGP with Monte Carlo Forward Sampling}} (Kaskade GP Eksak berantai), bukan \textit{Deep Gaussian Process} sejati.
\end{enumerate}

\vspace{0.2cm}

% =========================================================================
% SEKSI 2: ANALISIS ELBO
% =========================================================================
\section{Kajian Pertanyaan 2: Apakah Terdapat Penggunaan ELBO?}
\label{sec:analisis_elbo}

\noindent \textbf{Jawabannya tegas: SAMA SEKALI TIDAK ADA.}

\subsection{Landasan Teori: Peran ELBO pada DGP}
Pada DGP berbasis inferensi variasional, \textit{Evidence Lower Bound} (ELBO) berfungsi sebagai fungsi objektif tunggal yang diturunkan melalui pertidaksamaan Jensen:
\begin{equation}
    \log p(\mathbf{y} \mid \mathbf{x}) \ge \mathcal{L}_{\mathrm{ELBO}} = \sum_{n=1}^N \underbrace{\mathbb{E}_{q(f_n^{(L)})}\left[ \log p(y_n \mid f_n^{(L)}) \right]}_{\text{Data Likelihood Fit}} - \sum_{\ell=1}^L \underbrace{\mathrm{KL}\left( q(\mathbf{u}^{(\ell)}) \,\Vert\, p(\mathbf{u}^{(\ell)}) \right)}_{\text{Kompleksitas / Regularisasi Prior}}
\end{equation}
ELBO menyeimbangkan antara kecocokan model terhadap data observasi (suku pertama) dan penalti divergensi informasi terhadap prior (suku kedua).

\subsection{Fakta Implementasi pada Notebook \texttt{CO2\_1.ipynb}}
Pada notebook \texttt{CO2\_1.ipynb}:
\begin{itemize}
    \item Tidak ada deklarasi kelas fungsi tujuan GPyTorch seperti \texttt{gpytorch.mlls.VariationalELBO} maupun \texttt{gpytorch.mlls.ExactMarginalLogLikelihood}.
    \item Tidak ditemukan satupun sintaks pelatihan: \texttt{loss.backward()} dan \texttt{optimizer.step()}.
    \item Model secara langsung dipasang pada mode inferensi non-gradien:
    \begin{codebox}
    \small
    \begin{verbatim}
    model_gp.eval()
    lh_gp.eval()
    with torch.no_grad(), gpytorch.settings.fast_pred_var():
        pred_gp_dense = model_gp(dense_x)
    \end{verbatim}
    \end{codebox}
\end{itemize}
Karena tidak ada proses estimasi parameter variasional, nilai ELBO tidak pernah dihitung, tidak pernah diaproksimasi, dan tidak pernah dimaksimalkan.

\vspace{0.2cm}

% =========================================================================
% SEKSI 3: ANALISIS OPTIMASI HYPERPARAMETER
% =========================================================================
\section{Kajian Pertanyaan 3: Apakah Melakukan Optimasi Hyperparameter?}
\label{sec:analisis_optimasi}

\noindent \textbf{Jawabannya tegas: TIDAK ADA OPTIMASI HYPERPARAMETER.}

\subsection{Parameter Bersifat Pilihan Manual (\textit{Fixed/Hardcoded})}
Seluruh parameter kernel pada notebook tersebut ditentukan secara sewenang-wenang (\textit{ad-hoc manual assignment}) di awal inisialisasi:
\begin{itemize}
    \item \textbf{GP 1-Layer (Sel 8)}:
    \begin{equation}
        \ell = 0{,}20 \text{ tahun}, \quad \sigma_f^2 = 1{,}0, \quad \sigma_n^2 = 0{,}05
    \end{equation}
    \item \textbf{DGP 2-Layer (Sel 10)}:
    \begin{align}
        \text{Layer 1} &: \ell_1 = 0{,}30, \quad \sigma_{n1}^2 = 0{,}01 \nonumber \\
        \text{Layer 2} &: \ell_2 = 0{,}30, \quad \sigma_{n2}^2 = 0{,}03
    \end{align}
    \item \textbf{DGP 3-Layer (Sel 12)}:
    \begin{align}
        \text{Layer 1} &: \ell_1 = 0{,}30, \quad \sigma_{n1}^2 = 0{,}01 \nonumber \\
        \text{Layer 2} &: \ell_2 = 0{,}30, \quad \sigma_{n2}^2 = 0{,}01 \nonumber \\
        \text{Layer 3} &: \ell_3 = 0{,}30, \quad \sigma_{n3}^2 = 0{,}03
    \end{align}
\end{itemize}

\subsection{Implikasi Matematis Ketiadaan Optimasi}
Dalam kaidah pembelajaran mesin berbasis Gaussian Process (Rasmussen \& Williams, 2006), nilai hyperparameter optimal $\hat{\boldsymbol{\theta}} = \{\hat{\ell}, \hat{\sigma}_f^2, \hat{\sigma}_n^2\}$ harus diperoleh melalui minimisasi \textit{Negative Log Marginal Likelihood} (NLML):
\begin{equation}
    \hat{\boldsymbol{\theta}} = \arg\min_{\boldsymbol{\theta}} \left\{ \frac{1}{2} \mathbf{y}^\top \mathbf{K}_y^{-1} \mathbf{y} + \frac{1}{2} \log \det \mathbf{K}_y + \frac{N}{2} \log 2\pi \right\}
\end{equation}
Pada notebook \texttt{CO2\_1.ipynb}, turunan berarah $\frac{\partial \mathcal{L}}{\partial \boldsymbol{\theta}}$ tidak pernah dihitung. Kurva prediksi yang terbentuk pada grafik semata-mata adalah hasil evaluasi rumus kondisional analitik GP standar:
\begin{equation}
    \mathbf{f}_* \mid \mathbf{y} \sim \mathcal{N}\left( K_{*N} (K_{NN} + \sigma_n^2 I)^{-1} \mathbf{y}, \; K_{**} - K_{*N}(K_{NN} + \sigma_n^2 I)^{-1} K_{N*} \right)
\end{equation}
dengan matriks kovariansi yang dihitung menggunakan angka hyperparameter tebakan manual.

\vspace{0.2cm}

% =========================================================================
% SEKSI 4: TABEL KOMPARASI MENYELURUH
% =========================================================================
\section{Tabel Komparasi Paradigma Pemodelan}
\label{sec:tabel_komparasi}

\begin{table}[H]
\centering
\small
\renewcommand{\arraystretch}{1.25}
\begin{tabularx}{\textwidth}{p{3.2cm} X X}
\toprule
\textbf{Aspek Teoretis} & \textbf{Implementasi \texttt{CO2\_1.ipynb}} & \textbf{DGP Modern DSVI Sejati (Salimbeni, 2017)} \\
\midrule
\textbf{Kelas Arsitektur} & Kaskade \texttt{ExactGP} deterministik & Hierarki Variasional \texttt{ApproximateGP} \\
\textbf{Basis Titik Latih} & Seluruh data training ($N = 1.289$) & Inducing points bebas ($M \ll N$, misal $M = 64$) \\
\textbf{Kompleksitas Memori} & $\mathcal{O}(N^2)$ Cholesky matriks penuh & $\mathcal{O}(M^2)$ per batch data \\
\textbf{Target Hidden Layer} & Dipaksa bernilai pasti ($h^{(1)} = \text{standarisasi}(x)$) & Variabel acak laten sejati ($h \sim q(h)$) \\
\textbf{Fungsi Objektif} & Tidak ada (tanpa loss function) & $\mathcal{L}_{\mathrm{ELBO}}$ (Likelihood fit $-$ penalti KL) \\
\textbf{Status Optimasi} & Tanpa optimasi (\textit{hardcoded fixed params}) & Optimasi gradien stokastik (Adam optimizer) \\
\textbf{Propagasi Laten} & Monte Carlo sampling saat evaluasi testing & Reparameterisasi pathwise saat training \& testing \\
\bottomrule
\end{tabularx}
\caption{Perbandingan struktural antara model pada \texttt{CO2\_1.ipynb} dan DGP DSVI formal.}
\label{tab:komparasi}
\end{table}

\vspace{0.2cm}

% =========================================================================
% SEKSI 5: PANDUAN STRATEGIS MAHASISWA MATEMATIKA
% =========================================================================
\section{Rekomendasi \& Panduan Tanggapan untuk Bimbingan}
\label{sec:rekomendasi}

\noindent Sebagai mahasiswa Departemen Matematika, kejujuran dan ketepatan formal notasi adalah nilai tertinggi dalam penulisan Tugas Akhir. Berikut arahan posisi yang harus Anda ambil:

\begin{alertbox}
\textbf{Bagaimana memposisikan file \texttt{CO2\_1.ipynb} di hadapan Dosen Pembimbing?}
\begin{enumerate}\setlength{\itemsep}{2pt}
    \item \textbf{Jangan klaim file tersebut sebagai DGP DSVI}: Jika Anda menyatakan notebook ini menggunakan DSVI atau mengoptimasi ELBO, dosen pembimbing akan segera melihat bahwa kodenya tidak memuat fungsi loss dan inducing points, sehingga kredibilitas teori Anda akan dipertanyakan.
    \item \textbf{Posisikan sebagai Simulasi Awal Perambatan Variansi (\textit{Toy Demonstration})}: Sampaikan bahwa eksperimen pada \texttt{CO2\_1.ipynb} adalah \textit{studi simulasi perilaku perambatan variansi kaskade GP} untuk membuktikan secara visual bagaimana pita variansi $\pm 2\sigma$ melebar ketika beberapa lapisan fungsi dievaluasi secara berantai via integrasi Monte Carlo.
    \item \textbf{Tunjukkan Arsitektur DSVI Sejati yang Sudah Anda Siapkan}: Arahkan dosen pembimbing bahwa implementasi DGP berbasis DSVI dan optimasi ELBO yang sesungguhnya telah Anda bangun secara sistematis pada folder repositori:
    \begin{center}
        \texttt{02\_eksperimen/01\_dgp\_gpytorch/01\_dgp\_gpytorch\_implementation.ipynb}
    \end{center}
    di mana pada notebook tersebut: \texttt{ApproximateGP}, \texttt{VariationalELBO}, $M$ buah titik induksi, serta pembaruan gradien menggunakan \texttt{torch.optim.Adam} telah bekerja secara utuh dan terverifikasi secara matematis.
\end{enumerate}
\end{alertbox}

\end{document}
